
\documentclass[11pt]{article}

\usepackage[a4paper,margin=27mm]{geometry}
\usepackage{amsmath,amssymb,amsthm,mathtools}
\usepackage{booktabs,tabularx,longtable,array}
\usepackage[hidelinks]{hyperref}
\usepackage{microtype}
\usepackage{xcolor}

\definecolor{wrehblue}{RGB}{24,86,128}

\newtheorem{definition}{Definition}[section]
\newtheorem{proposition}[definition]{Proposition}
\newtheorem{example}[definition]{Example}
\newtheorem{remark}[definition]{Remark}

\newcommand{\W}{\mathcal W}
\newcommand{\Pcal}{\mathcal P}
\newcommand{\Acal}{\mathcal A}
\newcommand{\Bcal}{\mathcal B}
\newcommand{\QA}{\mathcal Q_{\mathcal A}}
\newcommand{\RA}{R_{\mathcal A}}

\title{\textbf{WR-I Mathematical Spine v0.1}\\[4pt]
\large Response-Defined World Spaces:\\
Finite-Resolution Equivalence and Identifiability of Global Realizations}
\author{A. A. Malachevsky\\
ORCID: 0009-0008-6009-3196\\
World Realizability \& Epistemic Horizons (WREH)\\
Research Programme \& Community}
\date{6 October 2026}

\begin{document}
\maketitle

\begin{abstract}
This document fixes the mathematical spine for the first technical paper of the
World Realizability \& Epistemic Horizons (WREH) programme.  The basic object is
a class of admissible global realizations equipped with a family of admissible
observation protocols and their response maps.  Exact response equality induces
a quotient world space.  We prove the canonical quotient representation,
factorization of identifiable quantities, monotonicity under protocol
refinement, factorization through gauge orbits, and a criterion for when
response non-uniqueness is exhausted by gauge redundancy.  We also separate
exact equivalence from finite-tolerance indistinguishability, for which
transitivity generally fails.  Smooth and infinite-dimensional examples show
how refinement can remove genuine non-identifiability or leave persistent
operational ambiguity.  The framework is deliberately ontologically neutral:
response-equivalent realizations are not declared physically identical, and no
claim is made that observation creates, selects, or modifies a world.
\end{abstract}

\section{Scope and claim ceiling}

WR-I is an abstract response-identifiability paper.  The term ``world'' means
an element of a declared class of global realizations.  It does not, by itself,
mean a cosmological universe or a physically realized ontology.

The paper does \emph{not} claim that:
\begin{itemize}
  \item response-equivalent realizations are ontologically identical;
  \item underlying reality does not exist;
  \item measurement creates or selects a unique world;
  \item finite-dimensional response walls are spacetime boundaries;
  \item protocol refinement is physical time, dynamics, or renormalization flow.
\end{itemize}

\section{World-response structures}

\begin{definition}[World-response structure]
A \emph{world-response structure} is a triple
\[
\mathbf W=(\W,\Pcal,\{R_P\}_{P\in\Pcal}),
\]
where $\W$ is a set of declared admissible global realizations, $\Pcal$ is a set
of admissible observation protocols, and each protocol $P$ has a response map
\[
R_P:\W\to Y_P
\]
into a protocol-dependent response space $Y_P$.
\end{definition}

\begin{definition}[Protocol bundle and joint response]
For $\Acal\subseteq\Pcal$, define
\[
Y_{\Acal}:=\prod_{P\in\Acal}Y_P
\]
and the joint response
\[
\RA:\W\to Y_{\Acal},
\qquad
\RA(W):=(R_P(W))_{P\in\Acal}.
\]
For the empty bundle, $Y_{\varnothing}$ is the one-point set.
\end{definition}

\begin{definition}[Exact response equivalence]
For a protocol bundle $\Acal$, define
\[
W_1\sim_{\Acal}W_2
\quad\Longleftrightarrow\quad
\RA(W_1)=\RA(W_2).
\]
Equivalently,
\[
R_P(W_1)=R_P(W_2)
\quad\text{for every }P\in\Acal.
\]
\end{definition}

\begin{definition}[Response-defined world space]
The \emph{response-defined world space} for $\Acal$ is the quotient set
\[
\QA:=\W/\!\sim_{\Acal}.
\]
Its canonical quotient map is denoted
\[
q_{\Acal}:\W\to\QA,
\qquad
q_{\Acal}(W)=[W]_{\Acal}.
\]
\end{definition}

\section{Canonical quotient representation}

\begin{proposition}[Response quotient representation]
The relation $\sim_{\Acal}$ is an equivalence relation.  Moreover, there is a
unique bijection
\[
\overline R_{\Acal}:\QA\longrightarrow \operatorname{Im}\RA
\]
such that
\[
\RA=\overline R_{\Acal}\circ q_{\Acal}.
\]
\end{proposition}

\begin{proof}
Equality in $Y_{\Acal}$ is reflexive, symmetric, and transitive.  Since
$\sim_{\Acal}$ is the pullback of equality along $\RA$, it is an equivalence
relation.

Define
\[
\overline R_{\Acal}([W]_{\Acal}):=\RA(W).
\]
If $[W]_{\Acal}=[W']_{\Acal}$, then $W\sim_{\Acal}W'$, hence
$\RA(W)=\RA(W')$.  Therefore the definition is well posed.  The displayed
factorization follows immediately.

The map is surjective onto $\operatorname{Im}\RA$ by definition of the image.
If
\[
\overline R_{\Acal}([W]_{\Acal})
=
\overline R_{\Acal}([W']_{\Acal}),
\]
then $\RA(W)=\RA(W')$, so $W\sim_{\Acal}W'$ and
$[W]_{\Acal}=[W']_{\Acal}$.  Hence $\overline R_{\Acal}$ is injective.
Uniqueness follows because $q_{\Acal}$ is surjective.
\end{proof}

\begin{remark}
As a bare set, $\QA$ contains exactly the information present in the joint
response image.  Additional geometry, topology, measure, smooth structure, or
stratification on $\QA$ requires additional assumptions and is not automatic.
\end{remark}

\section{Identifiable quantities}

\begin{definition}[Identifiable quantity]
Let $F:\W\to Z$.  We say that $F$ is \emph{$\Acal$-identifiable} if
\[
W_1\sim_{\Acal}W_2
\quad\Longrightarrow\quad
F(W_1)=F(W_2).
\]
Thus $F$ is constant on each response fibre.
\end{definition}

\begin{proposition}[Identifiability factorization]
A map $F:\W\to Z$ is $\Acal$-identifiable if and only if there exists a unique
map
\[
\overline F:\QA\to Z
\]
such that
\[
F=\overline F\circ q_{\Acal}.
\]
\end{proposition}

\begin{proof}
Assume first that $F$ is $\Acal$-identifiable.  Define
\[
\overline F([W]_{\Acal}):=F(W).
\]
If $[W]_{\Acal}=[W']_{\Acal}$, identifiability gives $F(W)=F(W')$, so
$\overline F$ is well defined.  The factorization follows directly.

Conversely, if $F=\overline F\circ q_{\Acal}$ and
$W_1\sim_{\Acal}W_2$, then
$q_{\Acal}(W_1)=q_{\Acal}(W_2)$, hence $F(W_1)=F(W_2)$.
Uniqueness again follows from surjectivity of $q_{\Acal}$.
\end{proof}

\begin{remark}
This proposition gives the operational content of identifiability: every
quantity determined by the declared protocol family must descend to the
response quotient.
\end{remark}

\section{Protocol refinement}

\begin{definition}[Refinement order]
For protocol bundles $\Acal,\Bcal\subseteq\Pcal$, write
\[
\Acal\preceq\Bcal
\quad\Longleftrightarrow\quad
\Acal\subseteq\Bcal.
\]
\end{definition}

\begin{proposition}[Refinement monotonicity]
If $\Acal\subseteq\Bcal$, then
\[
W_1\sim_{\Bcal}W_2
\quad\Longrightarrow\quad
W_1\sim_{\Acal}W_2,
\]
and therefore
\[
[W]_{\Bcal}\subseteq[W]_{\Acal}
\qquad
\text{for every }W\in\W.
\]
There is a unique surjection
\[
\pi_{\Bcal\Acal}:\mathcal Q_{\Bcal}\to\mathcal Q_{\Acal},
\qquad
[W]_{\Bcal}\mapsto[W]_{\Acal},
\]
satisfying
\[
q_{\Acal}
=
\pi_{\Bcal\Acal}\circ q_{\Bcal}.
\]
For $\Acal\subseteq\Bcal\subseteq\mathcal C$,
\[
\pi_{\mathcal C\Acal}
=
\pi_{\Bcal\Acal}\circ\pi_{\mathcal C\Bcal}.
\]
\end{proposition}

\begin{proof}
Equality of all responses indexed by $\Bcal$ implies equality of the subfamily
indexed by $\Acal$.  Hence $\sim_{\Bcal}$ refines $\sim_{\Acal}$.

If $[W]_{\Bcal}=[W']_{\Bcal}$, then $W\sim_{\Bcal}W'$, and therefore
$W\sim_{\Acal}W'$.  Thus $\pi_{\Bcal\Acal}$ is well defined.  It is surjective
because every $\Acal$-class has a representative $W$, and the $\Bcal$-class of
that same representative maps to it.  The commuting identity and the
composition law are immediate from the definition.
\end{proof}

\begin{remark}
The family $(\mathcal Q_{\Acal},\pi_{\Bcal\Acal})$ is the natural input for the
global-realizability questions of WR-II.  WR-I does not assume that any inverse
limit has a physically meaningful interpretation.
\end{remark}

\section{Gauge redundancy and response non-identifiability}

\begin{definition}[Gauge action]
Let a group $G$ act on $\W$.  The action is \emph{response-invariant} for
$\Acal$ if
\[
R_P(g\cdot W)=R_P(W)
\]
for every $g\in G$, $W\in\W$, and $P\in\Acal$.
\end{definition}

\begin{proposition}[Gauge factorization]
If the $G$-action is response-invariant for $\Acal$, then
\[
G\cdot W\subseteq[W]_{\Acal}
\]
for every $W\in\W$.  Furthermore, $\RA$ factors uniquely through the orbit
quotient:
\[
\W\xrightarrow{\pi_G}\W/G
\xrightarrow{\widetilde R_{\Acal}}Y_{\Acal},
\]
where
\[
\widetilde R_{\Acal}(G\cdot W)=\RA(W).
\]
\end{proposition}

\begin{proof}
For every $g\in G$, response invariance gives
$\RA(g\cdot W)=\RA(W)$, so $g\cdot W\sim_{\Acal}W$.  Thus the entire orbit is
contained in the response class.

Define
\[
\widetilde R_{\Acal}(G\cdot W):=\RA(W).
\]
If $W'$ lies in the same orbit as $W$, response invariance gives
$\RA(W')=\RA(W)$; hence the map is well defined.  The factorization is
immediate, and uniqueness follows from surjectivity of $\pi_G$.
\end{proof}

\begin{proposition}[Gauge-completeness criterion]
Assume the $G$-action is response-invariant for $\Acal$.  Then the following are
equivalent:
\begin{enumerate}
  \item $[W]_{\Acal}=G\cdot W$ for every $W\in\W$;
  \item the induced map
  \[
  \widetilde R_{\Acal}:\W/G\to Y_{\Acal}
  \]
  is injective.
\end{enumerate}
\end{proposition}

\begin{proof}
Assume (1).  Suppose
\[
\widetilde R_{\Acal}(G\cdot W_1)
=
\widetilde R_{\Acal}(G\cdot W_2).
\]
Then $\RA(W_1)=\RA(W_2)$, so $W_1\sim_{\Acal}W_2$.  By (1),
$W_2\in[W_1]_{\Acal}=G\cdot W_1$.  Thus the two orbits are equal, proving
injectivity.

Conversely, assume (2).  Proposition 5.2 already gives
\[
G\cdot W\subseteq[W]_{\Acal}.
\]
Let $W'\in[W]_{\Acal}$.  Then $\RA(W')=\RA(W)$, so
\[
\widetilde R_{\Acal}(G\cdot W')
=
\widetilde R_{\Acal}(G\cdot W).
\]
Injectivity implies $G\cdot W'=G\cdot W$, hence $W'\in G\cdot W$.  Therefore
$[W]_{\Acal}\subseteq G\cdot W$, proving equality.
\end{proof}

\begin{remark}
The criterion separates representational redundancy from genuine response
non-identifiability.  If $\widetilde R_{\Acal}$ is non-injective, then distinct
gauge orbits remain observationally indistinguishable.
\end{remark}


\section{Finite-resolution separation on compact physical state spaces}

The previous propositions are set-theoretic.  To connect exact identifiability
with finite experimental design, assume now that the gauge quotient
\[
X:=\W/G
\]
is a compact metric space with metric $d_X$, that each gauge-invariant protocol
induces a continuous map
\[
\widetilde R_P:X\to Y_P
\]
into a metric space $(Y_P,d_P)$, and that the full admissible family separates
points of $X$:
\[
x\neq y
\quad\Longrightarrow\quad
\exists P\in\Pcal:
\widetilde R_P(x)\neq\widetilde R_P(y).
\]

\begin{proposition}[Finite-resolution separation theorem]
Under the assumptions above, for every $\delta>0$ there exist a finite protocol
bundle
\[
\Acal_\delta=\{P_1,\dots,P_m\}\subseteq\Pcal
\]
and a constant $\eta_\delta>0$ such that
\[
d_X(x,y)\ge\delta
\]
implies
\[
\max_{1\le j\le m}
d_{P_j}\!\left(
\widetilde R_{P_j}(x),
\widetilde R_{P_j}(y)
\right)
\ge\eta_\delta.
\]
Thus every pair of physical states separated by at least $\delta$ in the
declared metric is distinguished by one of finitely many protocols, with a
uniform positive response margin.
\end{proposition}

\begin{proof}
Define the compact subset
\[
K_\delta
=
\{(x,y)\in X\times X:d_X(x,y)\ge\delta\}.
\]
For each protocol $P$, let
\[
U_P
=
\{(x,y)\in K_\delta:
d_P(\widetilde R_P(x),\widetilde R_P(y))>0\}.
\]
Continuity of $\widetilde R_P$ implies that $U_P$ is open in $K_\delta$.
Because the full protocol family separates points and every pair in
$K_\delta$ is distinct, the family $\{U_P\}_{P\in\Pcal}$ covers $K_\delta$.
Compactness yields a finite subcover
\[
U_{P_1},\dots,U_{P_m}.
\]

Now define
\[
h(x,y)
=
\max_{1\le j\le m}
d_{P_j}(\widetilde R_{P_j}(x),\widetilde R_{P_j}(y)).
\]
The function $h$ is continuous and strictly positive on $K_\delta$, because
the selected $U_{P_j}$ cover $K_\delta$.  By compactness,
\[
\eta_\delta
:=
\min_{(x,y)\in K_\delta}h(x,y)
>0.
\]
This is the required uniform separation margin.
\end{proof}

\begin{remark}
The proposition does not assert that one finite protocol bundle identifies all
points exactly.  As $\delta\downarrow0$, the required bundle size may grow and
the margin $\eta_\delta$ may tend to zero.  The result is therefore a
finite-resolution statement, not a finite exact-reconstruction theorem.
\end{remark}

\begin{remark}
Compactness, the chosen metric on $X$, continuity of the protocol responses,
and joint point separation are substantive assumptions.  In a physical
application each must be independently justified; none follows merely from the
word ``world''.
\end{remark}


\section{Finite tolerance is not an equivalence relation in general}

Assume each $Y_P$ carries a metric $d_P$, and choose positive weights $w_P$ for
which the following supremum is finite on the pairs under consideration:
\[
d_{\Acal}(W_1,W_2)
=
\sup_{P\in\Acal}
w_P\,d_P(R_P(W_1),R_P(W_2)).
\]

For $\varepsilon>0$, define the tolerance relation
\[
W_1\approx_{\Acal,\varepsilon}W_2
\quad\Longleftrightarrow\quad
d_{\Acal}(W_1,W_2)\le\varepsilon.
\]

\begin{proposition}[Failure of transitivity at finite tolerance]
The relation $\approx_{\Acal,\varepsilon}$ is generally not transitive and
therefore does not, in general, define a quotient.
\end{proposition}

\begin{proof}
Take one scalar response and three realizations with
\[
R(W_0)=0,\qquad
R(W_1)=\frac{3}{4}\varepsilon,\qquad
R(W_2)=\frac{3}{2}\varepsilon.
\]
Then
\[
|R(W_0)-R(W_1)|=\frac34\varepsilon\le\varepsilon,
\]
and
\[
|R(W_1)-R(W_2)|=\frac34\varepsilon\le\varepsilon,
\]
but
\[
|R(W_0)-R(W_2)|=\frac32\varepsilon>\varepsilon.
\]
Hence
\[
W_0\approx_{\varepsilon}W_1,\qquad
W_1\approx_{\varepsilon}W_2,\qquad
W_0\not\approx_{\varepsilon}W_2.
\]
\end{proof}

\begin{remark}
WR-I therefore reserves ``equivalence class'' for exact response equality unless
an explicit equivalence closure or another mathematically justified coarse
relation is introduced.  Finite-precision objects are called tolerance fibres,
response neighbourhoods, or uncertainty regions.
\end{remark}

\section{Examples}

\begin{example}[Finite binary world with gauge and hidden physical content]
Let
\[
\W=\{0,1\}^4,\qquad W=(a,b,g,h).
\]
Let $G=\mathbb Z_2$ act by flipping $g$ and leaving $(a,b,h)$ fixed.  Define
\[
R_a(W)=a,\qquad
R_b(W)=b,\qquad
R_h(W)=a\oplus h.
\]
With only $R_a$, there are two response classes of size eight.  With
$(R_a,R_b)$, there are four classes of size four.  With
$(R_a,R_b,R_h)$, there are eight classes of size two, and each class is exactly
one gauge orbit.  Thus the full admissible response map is injective on $\W/G$.

Adding an artificial protocol $R_g(W)=g$ destroys gauge invariance.  It
illustrates why the admissible protocol family must be part of the model rather
than an arbitrary list of mathematical readouts.
\end{example}

\begin{example}[Smooth world space: refinement from genuine ambiguity to pure gauge]
Let
\[
\W=S^1_x\times S^1_g,
\]
and let the gauge group $G=S^1$ act by translations in the $g$ coordinate:
\[
\alpha\cdot(x,g)=(x,g+\alpha).
\]
Then
\[
\W/G\cong S^1_x.
\]
Consider two gauge-invariant protocols
\[
R_c(x,g)=\cos x,
\qquad
R_s(x,g)=\sin x.
\]

For $\Acal=\{R_c\}$, the induced map on the gauge quotient is
\[
\widetilde R_{\Acal}:S^1\to[-1,1],
\qquad
x\mapsto\cos x.
\]
It is not injective.  For generic $x$, the response fibre contains two distinct
gauge orbits, represented by $x$ and $-x$.  Thus the ambiguity is not exhausted
by gauge.

For $\Bcal=\{R_c,R_s\}$,
\[
\widetilde R_{\Bcal}:S^1\to S^1\subset\mathbb R^2,
\qquad
x\mapsto(\cos x,\sin x)
\]
is injective.  Hence
\[
[W]_{\Bcal}=G\cdot W
\]
for every $W$.  Protocol refinement has removed genuine response
non-identifiability while leaving gauge redundancy intact.
\end{example}

\begin{example}[Continuous histories and dense observations]
Let
\[
\W=C([0,1],\mathbb R)
\]
be the space of continuous real-valued histories.  For each $t\in[0,1]$, define
the evaluation protocol
\[
P_t(f)=f(t).
\]

For a finite set
\[
\Acal=\{t_1,\dots,t_n\},
\]
the response fibre through $f$ is
\[
[f]_{\Acal}
=
\{g\in C([0,1]):g(t_i)=f(t_i)\text{ for }i=1,\dots,n\}.
\]
Equivalently,
\[
[f]_{\Acal}
=
f+V_{\Acal},
\qquad
V_{\Acal}
=
\{u\in C([0,1]):u(t_i)=0\text{ for all }i\}.
\]
The vector space $V_{\Acal}$ is infinite-dimensional, so every finite sampling
bundle leaves an infinite-dimensional ambiguity.

Now let $D\subset[0,1]$ be dense and countable.  If
\[
f(t)=g(t)\qquad\text{for every }t\in D,
\]
continuity implies $f=g$ on all of $[0,1]$.  Hence the protocol family
$\{P_t:t\in D\}$ is jointly injective on $\W$.

The conclusion depends essentially on the admissible realization class.  If
$\W$ were the set of all functions $[0,1]\to\mathbb R$ with no continuity
assumption, equality on a dense subset would not imply equality everywhere.
Thus identifiability is controlled jointly by the protocol family and the
admissibility conditions on $\W$.
\end{example}

\section{External benchmark: anisotropic Calderón gauge}

A standard benchmark is the anisotropic Calderón inverse problem.  The
Dirichlet-to-Neumann map for a Riemannian metric is invariant under
diffeomorphisms that fix the boundary.  Thus boundary-fixing diffeomorphism
orbits lie inside response fibres.  Rigidity results seek conditions under which
equality of the boundary response implies equality up to precisely this gauge.
This is a concrete domain-specific analogue of the gauge-completeness criterion.

WR-I does not claim to rediscover this phenomenon.  It uses it as a control case
for the abstract distinction between:
\[
\text{representational redundancy}
\qquad\text{and}\qquad
\text{genuine response non-identifiability}.
\]

\section{Prior-art boundary}

\begin{longtable}{>{\raggedright\arraybackslash}p{0.22\textwidth}
                  >{\raggedright\arraybackslash}p{0.34\textwidth}
                  >{\raggedright\arraybackslash}p{0.34\textwidth}}
\toprule
\textbf{Prior line} & \textbf{Established idea} & \textbf{WR-I boundary} \\
\midrule
\endfirsthead
\toprule
\textbf{Prior line} & \textbf{Established idea} & \textbf{WR-I boundary} \\
\midrule
\endhead

Rothenberg (1971), structural identification &
Observational equivalence and identifiability of structures in parametric
stochastic models; identification asks whether observable laws distinguish the
underlying structure. &
WR-I does not claim novelty for observational equivalence.  It packages
admissible \emph{global realizations}, a protocol-indexed response family,
gauge separation, and a refinement system intended to feed a later
global-realizability problem. \\[6pt]

Zhang (2026), Mathematical Identifiability Geometry &
Observation-equivalence classes are treated geometrically; indistinguishability
strata and unique-identification points are explicit geometric objects. &
WR-I must not claim that ``geometry of equivalence classes'' is new.  Its
distinctive target is a physically admissible world-response structure and the
separation of gauge, protocol refinement, and downstream global completion. \\[6pt]

Manchak (2009), global spacetime underdetermination &
Classical general relativity admits observationally indistinguishable
spacetimes with different global structure. &
Manchak is a domain-specific predecessor for global-world
underdetermination.  WR-I abstracts away from the GR-specific definition of
observational access and makes the admissible protocol family explicit. \\[6pt]

Cinti--Fano (2021), critique of Manchak &
Observationally indistinguishable spacetimes may fail proposed criteria of
physical reasonableness. &
This supports treating the admissible realization class
$\W_{\mathrm{admissible}}$ as part of the mathematical input rather than using
all formally constructible models. \\[6pt]

Abramsky--Brandenburger (2011), contextuality &
Compatible local empirical data may fail to admit a global section; global
realizability can be obstructed. &
This is a major predecessor for WR-II, not a novelty claim of WR-I.  WREH will
need to distinguish global outcome assignments from global world
realizations. \\[6pt]

Anisotropic Calderón inverse problem &
Boundary responses are invariant under boundary-fixing diffeomorphisms;
uniqueness is naturally formulated up to gauge. &
Provides a benchmark for Propositions 5.2--5.3 and for the distinction between
gauge redundancy and residual non-identifiability. \\
\bottomrule
\end{longtable}

\section{Novelty claim permitted at v0.1}

The defensible novelty claim at this stage is architectural rather than
elementary-theorem novelty:

\begin{quote}
WR-I formulates a protocol-indexed response framework whose primary state
space is a declared class of admissible global realizations, separates gauge
redundancy from residual response non-identifiability, organizes protocol
refinement by canonical quotient maps, and connects ideal joint
identifiability to finite experimental separation at any fixed nonzero
resolution on compact metric gauge quotients.  The resulting hierarchy is
then prepared for a distinct global-realizability problem.
\end{quote}

This claim must remain subject to further literature audit.

\section{Handoff to WR-II}

WR-II may ask whether a refinement family of nonempty finite-response fibres
admits a globally compatible realization.  In the language established here,
one may study systems such as
\[
\mathcal Q_1
\leftarrow
\mathcal Q_2
\leftarrow
\mathcal Q_3
\leftarrow\cdots
\]
and the compatibility or failure of global completion.

The present paper does not identify:
\[
\text{finite consistency}
\quad\text{with}\quad
\text{global existence}.
\]
Compactness assumptions, inverse/projective limits, sheaf-theoretic
obstructions, and explicit counterexamples belong to the next paper.

\begin{thebibliography}{99}

\bibitem{Rothenberg1971}
T. J. Rothenberg,
\emph{Identification in Parametric Models},
Econometrica 39 (1971), 577--591.
DOI: 10.2307/1913267.

\bibitem{Zhang2026}
J. Zhang,
\emph{Mathematical Identifiability Geometry},
Zenodo preprint (2026).
DOI: 10.5281/zenodo.22023251.

\bibitem{Manchak2009}
J. B. Manchak,
\emph{Can we know the global structure of spacetime?},
Studies in History and Philosophy of Modern Physics 40 (2009), 53--56.
DOI: 10.1016/j.shpsb.2008.07.004.

\bibitem{CintiFano2021}
E. Cinti and V. Fano,
\emph{Careful with those scissors, Eugene! Against the observational
indistinguishability of spacetimes},
Studies in History and Philosophy of Science 89 (2021), 103--113.
DOI: 10.1016/j.shpsa.2021.07.007.

\bibitem{AbramskyBrandenburger2011}
S. Abramsky and A. Brandenburger,
\emph{The Sheaf-Theoretic Structure of Non-Locality and Contextuality},
New Journal of Physics 13 (2011), 113036.
arXiv:1102.0264.

\bibitem{CalderonPoisson2019}
M. Lassas, T. Liimatainen, M. Salo,
\emph{The Poisson embedding approach to the Calderón problem},
Mathematische Annalen (published online 2019).
The Dirichlet-to-Neumann map is invariant under boundary-fixing
diffeomorphisms.

\end{thebibliography}

\end{document}
